\documentclass[11pt]{letter}

\usepackage[margin=1in]{geometry}
\usepackage[T1]{fontenc}
\usepackage{url}

\signature{Vincenzo Barone\\
INSTM, National Interuniversity Consortium for Materials Science and Technology\\
via G. Giusti 9\\
50121 Firenze, Italy\\
\texttt{vincebarone52@gmail.com}}

\address{Vincenzo Barone\\
INSTM, National Interuniversity Consortium for Materials Science and Technology\\
via G. Giusti 9\\
50121 Firenze, Italy\\
\texttt{vincebarone52@gmail.com}}

\date{\today}

\begin{document}

\begin{letter}{Editor\\
Journal of Chemical Theory and Computation}

\opening{Dear Editor,}

Please consider the manuscript entitled
``Non-Redundant and Symmetry-Adapted Coordinates for Black-Box Molecular
Structure Refinement'' for publication in the \emph{Journal of Chemical
Theory and Computation}.

The manuscript addresses a persistent limitation of semiexperimental molecular
structure refinement: the fitted geometry often depends on a user-defined
Z-matrix or on manually selected internal-coordinate models. This dependence is
benign for simple acyclic molecules, but it becomes problematic for rings,
fused systems, bridged frameworks, planar molecules, and cases with incomplete
isotopic substitution data.

We introduce \texttt{SEfit}, a framework in which the structure is read as
Cartesian coordinates and the active refinement space is generated
automatically from molecular topology and symmetry. Two non-redundant,
symmetry-adapted coordinate representations are available: generalized internal
coordinates (GICs) built from homogeneous primitive blocks, and
symmetry-adapted Cartesian displacements with
translations and rotations removed. Chemical constraints, predicates, and
parameter classes remain defined in primitive internal coordinates and are
projected onto the chosen active basis, so the method preserves chemical
interpretability without requiring a Z-matrix.

The work extends the semiexperimental structure-refinement capabilities
developed in the \texttt{MSR} molecular-structure-refinement workflow. It
retains moment-of-inertia fitting, vibrational and electronic corrections,
parameter classes, reduced-dimensionality hydrogen models, covariance
propagation, Kraitchman diagnostics, and human-readable reports, while moving
coordinate construction and constraint handling to a fully automatic
Cartesian/topology layer.

The approach is validated on glycolaldehyde, glycine II, cyclopentadiene,
nitrobenzene, azulene, and norcamphor. These examples span standard acyclic,
functional-constraint, cyclic, planar aromatic, fused-ring, and bridged cases.
The results show that the
symmetry-adapted GIC and Cartesian active spaces give consistent ranks,
effective dimensionalities, and rotational residuals, while final bond lengths,
angles, dihedrals, and their uncertainties are obtained from the optimized
Cartesian structures with rigorous error propagation.

We believe the manuscript is appropriate for the \emph{Journal of Chemical
Theory and Computation} (JCTC) because it presents a general computational
methodology for coordinate-dependent inverse problems, not only a
molecule-specific refinement. The same coordinate layer is designed to support
semiexperimental structures, Gaussian input generation, Wilson G/F matrix and
potential-energy-distribution (PED) analysis, and related vibrational
workflows.

The manuscript is original, has not been published previously, and is not under
consideration elsewhere. The author has approved the submission and declares no
competing financial interest. The software, benchmark inputs, generated tables,
and regression tests can be made available to reviewers and, upon publication,
under the distribution model selected by the author.

\closing{Sincerely,}

\end{letter}

\end{document}
